\begin{frame}[fragile]{실습: numpy로 2층 신경망 + XOR}
\begin{lstlisting}
import math, random

def sigmoid(x):
    return 1 / (1 + math.exp(-x))

def sigmoid_prime(x):
    s = sigmoid(x)
    return s * (1 - s)

def two_layer_forward(x, W1, b1, W2, b2):
    z1 = [sum(W1[i][j] * x[j] for j in range(len(x))) + b1[i]
          for i in range(len(b1))]
    a1 = [sigmoid(v) for v in z1]
    z2 = sum(W2[i] * a1[i] for i in range(len(a1))) + b2
    return sigmoid(z2), (x, z1, a1, z2, sigmoid(z2))

def two_layer_backward(y_true, cache, W2):
    x, z1, a1, z2, a2 = cache
    delta2 = (a2 - y_true) * sigmoid_prime(z2)
    delta1 = [W2[i] * delta2 * sigmoid_prime(z1[i]) for i in range(len(z1))]
    grad_W2 = [delta2 * a1[i] for i in range(len(a1))]
    grad_W1 = [[delta1[i] * x[j] for j in range(len(x))]
               for i in range(len(delta1))]
    return grad_W1, delta1, grad_W2, delta2
\end{lstlisting}
  {\footnotesize 은닉 유닛 4개 + 5,000 에폭, 무작위 초기화:
  최종 예측 $[0.016,\; 0.977,\; 0.976,\; 0.030]$ $\to$ 반올림
  $[0, 1, 1, 0]$ — 정답 XOR 패턴과 정확히 일치.
  로지스틱회귀 혼자서는 영원히 못 풀던 문제를 \textbf{은닉층 하나 +
  역전파}로 풀어냄.}
\end{frame}
